\documentclass[12pt,a4paper]{article}
\usepackage[top=22mm,bottom=24mm,left=27mm,right=27mm]{geometry}
\usepackage{amsmath,amssymb}
\usepackage{xcolor}
\definecolor{link}{HTML}{205973}
\usepackage[colorlinks=true,linkcolor=link,urlcolor=link]{hyperref}
\hypersetup{pdftitle={Discounted state occupancy and the Bellman residual}}
\setlength{\parindent}{0pt}
\setlength{\parskip}{5pt}
\raggedbottom
\newcommand{\R}{\mathbb R}
\newcommand{\lead}[1]{\par\smallskip\noindent\textbf{#1}\enspace}

\begin{document}
{\LARGE\bfseries Discounted state occupancy\\[2pt] and the Bellman residual\par}
\medskip
For a fixed policy, the error in an approximate cost-to-go can be recovered
exactly from its one-step errors. The weight on each state is the fraction
of discounted time the policy is expected to spend there. This distribution has a
balance equation that determines it uniquely.

\section*{A fixed policy and its one-step error}
Let \(\mathcal S\) and \(\mathcal A\) be finite, nonempty sets of states and
actions. At state \(x\), action \(a\) costs \(g_{x,a}\in\R\) and moves to
state \(y\) with probability \(p(x,y,a)\). Fix a policy \(\mu\) that chooses
one action \(\mu(x)\) in each state. Its cost vector and transition matrix are
\[
  g_\mu(x)=g_{x,\mu(x)},\qquad
  P_\mu(x,y)=p(x,y,\mu(x)),\qquad \sum_y P_\mu(x,y)=1.
\]
Throughout, \(\alpha\in(0,1)\) is the discount factor. Values are column
vectors. An initial distribution \(\nu\) has \(\nu_x\ge0\) and
\(\sum_x\nu_x=1\). We write \(\nu^\top\) for its row vector, so
\(\nu^\top f=\sum_x\nu_x f(x)\) is the expected value of \(f\) at the start.

Starting from \(x\), the expected cost at time \(t\) is \((P_\mu^t g_\mu)(x)\),
where \(P_\mu^0=I\), the identity matrix. For a proposed cost-to-go \(J\),
the true discounted cost \(J_\mu\) and the one-step update \(T_\mu J\) are
\[
 J_\mu=\sum_{t=0}^{\infty}\alpha^t P_\mu^t g_\mu,
 \qquad T_\mu J=g_\mu+\alpha P_\mu J .
\]
The Bellman residual \(J-T_\mu J\) is positive at a state when \(J\)
exceeds the cost of taking one policy step and then using \(J\) thereafter.

\section*{The occupancy and residual identities}
\textbf{Proposition.} For every \(\mu\) and \(\nu\) as above, the row vector
\begin{equation}
 \pi_{\mu,\nu}^{\top}
 =(1-\alpha)\sum_{t=0}^{\infty}\alpha^t\nu^\top P_\mu^t
 =(1-\alpha)\nu^\top(I-\alpha P_\mu)^{-1}
 \label{eq:occ}
\end{equation}
is a probability distribution. It is the unique real row vector satisfying
\begin{equation}
 \pi_{\mu,\nu}^{\top}=(1-\alpha)\nu^\top+\alpha\pi_{\mu,\nu}^{\top}P_\mu .
 \label{eq:balance}
\end{equation}
For every \(J\in\R^{\mathcal S}\), the expected initial value error satisfies
\begin{equation}
 \nu^\top(J-J_\mu)
   =\frac{1}{1-\alpha}\,\pi_{\mu,\nu}^{\top}(J-T_\mu J).
 \label{eq:residual}
\end{equation}

\newpage
\section*{Proof}
\lead{The distribution of discounted visits.}
The row vector \(\nu^\top P_\mu^t\) is the state distribution at time \(t\).
Indeed, multiplying a probability row by \(P_\mu\) preserves
nonnegativity and total mass. The coefficients \((1-\alpha)\alpha^t\)
are nonnegative and sum to one. Thus their mixture in \eqref{eq:occ}
converges coordinatewise and has total mass one. A state with zero initial
mass can have positive occupancy if the policy can reach it later.

\lead{The inverse and the balance equation.}
Every entry of \(P_\mu^t\) lies between zero and one, since each row is
a probability distribution. The matrix series \(\sum_{t\ge0}\alpha^tP_\mu^t\)
therefore converges absolutely in each entry. For the partial sum
\(D_n=\sum_{t=0}^{n}\alpha^tP_\mu^t\), multiplication on either side gives
\[
  (I-\alpha P_\mu)D_n
      =D_n(I-\alpha P_\mu)=I-\alpha^{n+1}P_\mu^{n+1}.
\]
The last term on the right tends to zero entrywise. There are only finitely
many states, so multiplication by \(I-\alpha P_\mu\) commutes with the
limit. The sum is consequently its two-sided inverse. This proves the
second formula in \eqref{eq:occ}.

Multiplying \eqref{eq:occ} on the right by \(I-\alpha P_\mu\) gives
\[
   \pi_{\mu,\nu}^\top(I-\alpha P_\mu)=(1-\alpha)\nu^\top.
\]
Rearranging gives \eqref{eq:balance}. Conversely, any real row vector
satisfying that balance equation must give the same product with
\(I-\alpha P_\mu\). Multiplication by the inverse determines the row
uniquely. In \eqref{eq:balance}, the first term accounts for time zero
and the second for visits after one or more transitions.

\lead{From a one-step error to the initial error.}
The series for \(J_\mu\) gives \(J_\mu=(I-\alpha P_\mu)^{-1}g_\mu\).
Hence \(J_\mu\) is the unique solution of \(J_\mu=g_\mu+\alpha P_\mu J_\mu\).
Subtracting this equality from the definition of the residual gives
\[
 (I-\alpha P_\mu)(J-J_\mu)
     =J-\alpha P_\mu J-g_\mu=J-T_\mu J .
\]
Multiply on the left by \(\pi_{\mu,\nu}^\top\) and use the balance equation.
The left side becomes \((1-\alpha)\nu^\top(J-J_\mu)\). Division by
\(1-\alpha>0\) gives \eqref{eq:residual}. \hfill\(\square\)

\newpage
\section*{A two-state example}
The states are \(s,t\), and the initial state is \(s\). From either state,
action \(A\) costs \(1\) and moves to \(s\) or \(t\) with equal probability.
Action \(B\) costs \(2\) and moves to \(t\) with certainty. The policy
takes \(A\) at \(s\) and \(B\) at \(t\). In state order \(s,t\),
\[
 P_\mu=\begin{pmatrix}1/2&1/2\\0&1\end{pmatrix},\qquad
 g_\mu=\begin{pmatrix}1\\2\end{pmatrix},\qquad
 \nu^\top=(1,0),\qquad \alpha=\tfrac12 .
\]
Writing \(u=\pi_{\mu,\nu}(s)\), the first coordinate of
\eqref{eq:balance} is \(u=1/2+u/4\). Thus the occupancy is
\((2/3,1/3)\). State \(t\) accounts for a third of discounted time even
though the initial distribution gives it zero mass.

At \(t\), the policy pays \(2\) forever, giving \(J_\mu(t)=4\). At \(s\),
the Bellman equation gives
\(J_\mu(s)=1+\tfrac14 J_\mu(s)+\tfrac14 J_\mu(t)\), so \(J_\mu(s)=8/3\).
Take the candidate \(J(s)=3,\ J(t)=2\). Its one-step predictions and
errors are
\[
 \begin{array}{c@{\qquad}r@{\quad}r@{\quad}r@{\quad}r}
   \text{State} & J_\mu & J & T_\mu J & J-T_\mu J\\ \hline
   s & 8/3 & 3 & 9/4 & \phantom{-}3/4\\
   t & 4   & 2 & 3   & -1
 \end{array}
\]
The residuals at the two states have opposite signs. Their contributions
partly cancel when averaged over discounted visits,
\[
 \pi_{\mu,\nu}^\top(J-T_\mu J)=\frac23\cdot\frac34+\frac13\cdot(-1)=\frac16 .
\]
Dividing by \(1-\alpha=1/2\) gives an initial error of \(1/3\). Direct
evaluation at the starting state gives the same answer,
\(J(s)-J_\mu(s)=3-8/3=1/3\).

\section*{Relation to the original result}
Equation (17) of Bhat, Farias, Moallemi and Zheng defines the sampling law
for an optimal policy \(\mu^*\). Taking \(\mu=\mu^*\) in \eqref{eq:occ}
gives their formula as printed. The proof applies to any fixed policy and
allows zeros in the initial distribution. Equation \eqref{eq:residual}
makes explicit how this sampling law weights signed one-step errors.

\smallskip
{\small
\textbf{Source.} N.~Bhat, V.~F.~Farias, C.~C.~Moallemi and A.~T.~Zheng,
\href{https://moallemi.com/ciamac/papers/kernel-adp-2012.pdf\#page=12}{\emph{Non-parametric Approximate Dynamic Programming via the Kernel Method}},
August 29, 2022, \S3.2.\par
\textbf{Formal proof.}
\href{https://github.com/nikhilpb/kadp/blob/main/Kadp/Occupancy.lean}{Lean occupancy and residual identities}
and \href{https://github.com/nikhilpb/kadp/blob/main/Kadp/Examples/Occupancy.lean}{two-state example}.
The series proof applies
\href{https://github.com/LeanEconomics/DynamicProgramming/blob/9d4783f3fb282444f8f3a91c4e6439293dd4bacc/FiniteStates/MarkovDecisionProcesses/MarkovDecisionProcesses/MDP.lean}{Robert Kirkby's policy-evaluation theorems}
to each state's indicator function.}
\end{document}
